\section{State-conditioned procedural retrieval}\label{sec:C16:trajectories}
Procedural retrieval should match the consumer's current situation rather than only the initial task description. Multi-Agent Transactive Memory indexes recent observation-action context as a key and a subsequent continuation as its value, enabling one agent to reuse another agent's experience~\cite{kim2026matm}. A learned reranker estimates downstream utility beyond initial similarity. The method concerns shared artifacts across heterogeneous agents; it does not require that the agents jointly train or communicate during each task.

The essential indexing boundary is temporal. A key must contain information available before the continuation begins; the value can contain the later actions offered as guidance. If a key accidentally includes the successful future result, retrieval can exploit information unavailable to a real consumer. This is a procedural form of leakage. Dataset construction should validate the cutoff at each indexed step and preserve environment boundaries when sliding windows are formed.

The reported ALFWorld comparison improves from 47.08\% without memory to 64.31\% with retrieval and learned reranking, while WebArena gains are smaller~\cite{kim2026matm}. These are author-reported population results under the paper's allocations and environments, not a universal memory-system ranking. Generating the experience repository and training rerankers also consume resources. Consumer-side savings should not be described as total lifecycle savings unless those production costs are included or explicitly amortized.

An executable continuation is not a command to replay blindly. A different consumer may have a different tool set, current object state, or authorization. The retrieved trajectory should be treated as evidence for a candidate plan whose preconditions must be checked. A replay mechanism that substitutes identifiers from the current task can still preserve a hidden assumption from the producer's environment. The interface should expose those assumptions where possible and allow abstention or fresh exploration.

HyperSkill's subtask and task routes provide a contrasting procedural index~\cite{xu2026hyperskill}. Rather than supplying only a raw continuation, they retrieve lessons and skills associated with relevant trajectories. This can improve abstraction and reduce raw context, while introducing extraction and consolidation error. MATM and HyperSkill therefore compare along representation, timing, feedback, and cost axes. An aggregate success difference without matched consumer and producer budgets would not isolate whether raw or abstract experience is the better transferable object.

The proposed evaluation separates near-duplicate reuse, compositional transfer, and cross-environment transfer. Near-duplicate retrieval tests useful caching but is weak evidence for general procedural learning. Compositional transfer holds out combinations of known subtasks. Cross-environment transfer changes tools or state dynamics and tests compatibility. Report all three rather than pooling them into one ``generalization'' score.

